\subsection{In-Place Memory Mutation}

The previous subsection described set calculations that create new sets. For example, \texttt{a | b} produces a result without changing \texttt{a} or \texttt{b}. Python also provides mutating versions of these operations. These methods change the left-hand set object in place.

The word ``destructive'' does not mean that the object becomes invalid. It means that the previous membership state is overwritten by a new membership state.

\begin{quote}
A non-mutating set calculation produces a new set. An in-place set update changes the existing set object.
\end{quote}

This distinction can be seen with \texttt{id()}.

\begin{verbatim}
names = {"Homer", "Marge"}
extra = {"Lisa", "Bart"}

print(f"before: {names!r}")
print(f"id before: {id(names)}")

names.update(extra)

print()
print(f"after:  {names!r}")
print(f"id after:  {id(names)}")
\end{verbatim}

Possible output:

\begin{verbatim}
before: {'Homer', 'Marge'}
id before: 2199812636800

after:  {'Bart', 'Lisa', 'Homer', 'Marge'}
id after:  2199812636800
\end{verbatim}

The set object has the same identity before and after the call. Its internal membership domain changed.

\subsubsection{Destructive Union Updates via \texttt{.update()}}

The method \texttt{.update()} performs an in-place union. It adds all elements from another iterable into the existing set.

\begin{verbatim}
left = {"spam", "eggs"}
right = {"eggs", "larch", 42}

print(f"left before: {left!r}")

result = left.update(right)

print(f"left after:  {left!r}")
print(f"right:       {right!r}")
print(f"result:      {result!r}")
\end{verbatim}

Possible output:

\begin{verbatim}
left before: {'eggs', 'spam'}
left after:  {42, 'eggs', 'larch', 'spam'}
right:       {42, 'eggs', 'larch'}
result:      None
\end{verbatim}

The method returns \texttt{None} because the useful result is the changed state of \texttt{left}. This is the same convention seen in other in-place mutation methods.

The non-mutating version and mutating version differ like this:

\begin{verbatim}
a = {"Arthur", "Brian"}
b = {"Brian", "Lancelot"}

new_set = a | b

print("after a | b")
print(f"a:       {a!r}")
print(f"b:       {b!r}")
print(f"new_set: {new_set!r}")

a.update(b)

print()
print("after a.update(b)")
print(f"a:       {a!r}")
print(f"b:       {b!r}")
\end{verbatim}

Possible output:

\begin{verbatim}
after a | b
a:       {'Arthur', 'Brian'}
b:       {'Lancelot', 'Brian'}
new_set: {'Arthur', 'Lancelot', 'Brian'}

after a.update(b)
a:       {'Arthur', 'Lancelot', 'Brian'}
b:       {'Lancelot', 'Brian'}
\end{verbatim}

The expression \texttt{a | b} leaves \texttt{a} unchanged. The call \texttt{a.update(b)} changes \texttt{a}.

The method can accept more than one iterable.

\begin{verbatim}
items = {"spam"}

items.update(["eggs", "spam"], ("larch", 42), {"ni"})

print(f"items: {items!r}")
print(f"type(items): {type(items)}")
\end{verbatim}

Possible output:

\begin{verbatim}
items: {'eggs', 42, 'spam', 'ni', 'larch'}
type(items): <class 'set'>
\end{verbatim}

Each source is traversed, and each hashable value is inserted into the target set. Duplicate values collapse automatically.

The inserted values must still be hashable.

\begin{verbatim}
items = {"spam"}

try:
    items.update(["eggs", ["bad", "list"]])
except TypeError as err:
    print(f"update failed: {err}")

print(f"items after failure: {items!r}")
\end{verbatim}

Possible output:

\begin{verbatim}
update failed: unhashable type: 'list'
items after failure: {'eggs', 'spam'}
\end{verbatim}

The set may already have been partly changed before the error occurred. In-place mutation is immediate; it is not a guaranteed all-or-nothing transaction.

\subsubsection{Destructive Intersection Updates via \texttt{.intersection\_update()}}

The method \texttt{.intersection\_update()} keeps only elements that also appear in the supplied iterable or iterables. It changes the target set in place.

\begin{verbatim}
allowed = {"Homer", "Marge", "Lisa", "Bart"}
present = {"Lisa", "Bart", "Burns"}

print(f"allowed before: {allowed!r}")

result = allowed.intersection_update(present)

print(f"allowed after:  {allowed!r}")
print(f"present:        {present!r}")
print(f"result:         {result!r}")
\end{verbatim}

Possible output:

\begin{verbatim}
allowed before: {'Lisa', 'Homer', 'Bart', 'Marge'}
allowed after:  {'Lisa', 'Bart'}
present:        {'Lisa', 'Burns', 'Bart'}
result:         None
\end{verbatim}

The set \texttt{allowed} was reduced. Only values also found in \texttt{present} survived.

The non-mutating and mutating forms differ like this:

\begin{verbatim}
a = {"spam", "eggs", "larch"}
b = {"eggs", "beans", "larch"}

new_set = a & b

print("after a & b")
print(f"a:       {a!r}")
print(f"b:       {b!r}")
print(f"new_set: {new_set!r}")

a.intersection_update(b)

print()
print("after a.intersection_update(b)")
print(f"a:       {a!r}")
print(f"b:       {b!r}")
\end{verbatim}

Possible output:

\begin{verbatim}
after a & b
a:       {'eggs', 'spam', 'larch'}
b:       {'eggs', 'beans', 'larch'}
new_set: {'eggs', 'larch'}

after a.intersection_update(b)
a:       {'eggs', 'larch'}
b:       {'eggs', 'beans', 'larch'}
\end{verbatim}

The operator \texttt{\&} creates a new result. The method \texttt{.intersection\_update()} overwrites the membership state of \texttt{a}.

The method can also accept several iterables. The target keeps only values found in all of them.

\begin{verbatim}
names = {"Jerry", "Elaine", "George", "Kramer"}
group1 = ["Jerry", "Elaine", "Newman"]
group2 = {"Elaine", "Jerry", "Puddy"}

names.intersection_update(group1, group2)

print(f"names: {names!r}")
\end{verbatim}

Possible output:

\begin{verbatim}
names: {'Elaine', 'Jerry'}
\end{verbatim}

The final set contains only names present in the original \texttt{names}, in \texttt{group1}, and in \texttt{group2}.

\subsubsection{Destructive Difference Updates via \texttt{.difference\_update()}}

The method \texttt{.difference\_update()} removes from the target set all elements found in the supplied iterable or iterables.

\begin{verbatim}
items = {"spam", "eggs", "larch", "beans"}
blocked = {"eggs", "beans"}

print(f"items before: {items!r}")

result = items.difference_update(blocked)

print(f"items after:  {items!r}")
print(f"blocked:      {blocked!r}")
print(f"result:       {result!r}")
\end{verbatim}

Possible output:

\begin{verbatim}
items before: {'eggs', 'beans', 'larch', 'spam'}
items after:  {'larch', 'spam'}
blocked:      {'eggs', 'beans'}
result:       None
\end{verbatim}

The target set lost the values that were present in \texttt{blocked}. The \texttt{blocked} set itself was not changed.

The operation is directional. It removes from the target. It does not compute a balanced disagreement between both sets.

\begin{verbatim}
a = {"Arthur", "Brian", "Lancelot"}
b = {"Brian", "Galahad"}

a.difference_update(b)

print(f"a after difference_update: {a!r}")
print(f"b:                         {b!r}")
\end{verbatim}

Possible output:

\begin{verbatim}
a after difference_update: {'Arthur', 'Lancelot'}
b:                         {'Brian', 'Galahad'}
\end{verbatim}

Only \texttt{"Brian"} was removed from \texttt{a}. The value \texttt{"Galahad"} did not matter because it was not present in \texttt{a}.

The method can receive more than one iterable.

\begin{verbatim}
names = {"Homer", "Marge", "Lisa", "Bart", "Maggie"}

names.difference_update(["Homer", "Bart"], {"Maggie"})

print(f"names: {names!r}")
\end{verbatim}

Possible output:

\begin{verbatim}
names: {'Lisa', 'Marge'}
\end{verbatim}

The target set removed all values found in either removal source.

The non-mutating version is \texttt{.difference()} or the \texttt{-} operator.

\begin{verbatim}
original = {"spam", "eggs", "larch"}
blocked = {"eggs"}

new_set = original - blocked

print(f"original after '-': {original!r}")
print(f"new_set:            {new_set!r}")

original.difference_update(blocked)

print()
print(f"original after difference_update: {original!r}")
\end{verbatim}

Possible output:

\begin{verbatim}
original after '-': {'eggs', 'spam', 'larch'}
new_set:            {'spam', 'larch'}

original after difference_update: {'spam', 'larch'}
\end{verbatim}

The first calculation creates a new set. The update method changes the existing set.

\subsubsection{Destructive Symmetric Difference Updates via \texttt{.symmetric\_difference\_update()}}

The method \texttt{.symmetric\_difference\_update()} changes the target set so that it keeps only values that appear in exactly one of the two sources: the original target or the supplied iterable.

\begin{verbatim}
left = {"spam", "eggs", "larch"}
right = {"eggs", "beans", "larch", 42}

print(f"left before: {left!r}")

result = left.symmetric_difference_update(right)

print(f"left after:  {left!r}")
print(f"right:       {right!r}")
print(f"result:      {result!r}")
\end{verbatim}

Possible output:

\begin{verbatim}
left before: {'eggs', 'larch', 'spam'}
left after:  {42, 'beans', 'spam'}
right:       {42, 'eggs', 'beans', 'larch'}
result:      None
\end{verbatim}

The values \texttt{"eggs"} and \texttt{"larch"} were removed because they appeared in both. The values \texttt{"beans"} and \texttt{42} were added because they appeared only in \texttt{right}. The value \texttt{"spam"} remained because it appeared only in the original \texttt{left}.

This operation can be understood as two simultaneous actions:

\begin{itemize}
    \item remove elements shared by both sides;
    \item add elements that appear only in the supplied iterable.
\end{itemize}

The non-mutating and mutating forms differ like this:

\begin{verbatim}
a = {"Jerry", "Elaine", "George"}
b = {"Jerry", "Elaine", "Kramer"}

new_set = a ^ b

print("after a ^ b")
print(f"a:       {a!r}")
print(f"b:       {b!r}")
print(f"new_set: {new_set!r}")

a.symmetric_difference_update(b)

print()
print("after a.symmetric_difference_update(b)")
print(f"a:       {a!r}")
print(f"b:       {b!r}")
\end{verbatim}

Possible output:

\begin{verbatim}
after a ^ b
a:       {'Elaine', 'George', 'Jerry'}
b:       {'Elaine', 'Kramer', 'Jerry'}
new_set: {'George', 'Kramer'}

after a.symmetric_difference_update(b)
a:       {'George', 'Kramer'}
b:       {'Elaine', 'Kramer', 'Jerry'}
\end{verbatim}

The operator \texttt{\textasciicircum{}} produces a new set. The method \texttt{.symmetric\_difference\_update()} rewrites the target set.

Unlike \texttt{.update()}, \texttt{.intersection\_update()}, and \texttt{.difference\_update()}, this method takes one main iterable argument.

\begin{verbatim}
base = {"Arthur", "Brian", "Lancelot"}
other = ["Brian", "Galahad", "Arthur"]

base.symmetric_difference_update(other)

print(f"base: {base!r}")
\end{verbatim}

Possible output:

\begin{verbatim}
base: {'Galahad', 'Lancelot'}
\end{verbatim}

The list was accepted as an iterable source. The result remains a set.

The in-place set update methods can be summarized as follows:

\begin{center}
\begin{tabular}{lll}
\textbf{Method} & \textbf{Meaning} & \textbf{Target Set After Call} \\
\hline
\texttt{.update()} & Destructive union & Everything from target and source \\
\texttt{.intersection\_update()} & Destructive intersection & Only values common to all sources \\
\texttt{.difference\_update()} & Destructive difference & Target values not found in sources \\
\texttt{.symmetric\_difference\_update()} & Destructive symmetric difference & Values found on exactly one side \\
\end{tabular}
\end{center}

All these methods return \texttt{None}. Their purpose is to mutate the existing set object.

\begin{verbatim}
items = {"spam", "eggs"}
other = {"eggs", "larch"}

print(f"start id: {id(items)}")
print(f"start items: {items!r}")

items.update(other)

print()
print(f"after update id: {id(items)}")
print(f"after update items: {items!r}")
\end{verbatim}

Possible output:

\begin{verbatim}
start id: 2199812636800
start items: {'eggs', 'spam'}

after update id: 2199812636800
after update items: {'eggs', 'larch', 'spam'}
\end{verbatim}

The object identity is stable. The membership domain changed in place.